\Askhsh{26605_SOLUTION}{1}
\begin{ylh}{1}
    \item [Ύλη από εικόνα]
\end{ylh}
\begin{alist}
\item [Α)]
\begin{rlist}
\item Ισχύει ότι $f^2(x) - 5 = x^2$ για κάθε $x \in \mathbb{R}$ ή $f^2(x) = x^2+5$ για κάθε $x \in \mathbb{R}$.
    $f(x) = 0 \Leftrightarrow f^2(x) = 0 \Leftrightarrow x^2+5=0$, αδύνατο. Οπότε $f(x) \ne 0$ για κάθε $x \in \mathbb{R}$.
\item Η συνάρτηση $f$ είναι συνεχής στο $\mathbb{R}$ με $f(x) \ne 0$ για κάθε $x \in \mathbb{R}$. Οπότε η $f$ διατηρεί πρόσημο στο $\mathbb{R}$. Δίνεται ότι $f(2)=3>0$, οπότε η συνάρτηση $f$ παίρνει μόνο θετικές τιμές για κάθε $x \in \mathbb{R}$. Ισχύει $f^2(x)=x^2+5 \Leftrightarrow |f(x)|=\sqrt{x^2+5}$, και επειδή η $f$ παίρνει μόνο θετικές τιμές για κάθε $x \in \mathbb{R}$, θα ισχύει $f(x)=\sqrt{x^2+5}$ για κάθε $x \in \mathbb{R}$.
\end{rlist}
\item [Β)]
\begin{rlist}
\item Αν $g(x)=x^2-\sigma\upsilon\nu x$, με $x \in \mathbb{R}$, $g'(x)=2x+\eta\mu x$ και $g''(x)=2+\sigma\upsilon\nu x$ για κάθε $x \in \mathbb{R}$.
    Παρατηρούμε ότι $g''(x)>0$ για κάθε $x \in \mathbb{R}$, αφού $1 \le 2+\sigma\upsilon\nu x \le 3$, και η συνάρτηση $g'(x)$ είναι συνεχής στο $\mathbb{R}$, οπότε η συνάρτηση $g'$ είναι γνησίως αύξουσα στο $\mathbb{R}$.
    Για $x < 0$ ισχύει $g'(x) < g'(0) = 0$, αφού η συνάρτηση $g'$ είναι γνησίως αύξουσα, ενώ για $x>0$ ισχύει $g'(x) > g'(0) = 0$. Άρα για τη συνάρτηση $g$ έχουμε:
    $g$ συνεχής στο $(-\infty, 0]$ με $g'(x) < 0$ στο $(-\infty, 0)$, άρα η συνάρτηση $g$ είναι γνησίως φθίνουσα στο $(-\infty, 0]$. Αντίστοιχα
    $g$ συνεχής στο $[0, +\infty)$ με $g'(x) > 0$ στο $(0, +\infty)$, άρα η συνάρτηση $g$ είναι γνησίως αύξουσα στο $[0, +\infty)$.
    (Η συνάρτηση $g$ παρουσιάζει ολικό ελάχιστο στο $0$ το $g(0)=-1$).
\item Η εξίσωση $f^2(x)=5+\sigma\upsilon\nu x$ με $x \in \mathbb{R}$, γράφεται ισοδύναμα $x^2+5=5+\sigma\upsilon\nu x \Leftrightarrow x^2-\sigma\upsilon\nu x=0 \Leftrightarrow g(x)=0$ με $x \in \mathbb{R}$. Ζητείται να δείξουμε ότι η εξίσωση $g(x)=0$ έχει δύο ρίζες αντίθετες στο $(-\pi, \pi)$ και δεν έχει άλλες ρίζες στο $\mathbb{R}$.

    Η συνάρτηση $g$ είναι συνεχής στο $[0,\pi]$, με
    $g(\pi)=\pi^2-\sigma\upsilon\nu \pi = \pi^2+1 > 0$
    $g(0)=-\sigma\upsilon\nu 0 = -1 < 0$
    Η συνάρτηση $g$ ικανοποιεί τις προϋποθέσεις του θεωρήματος Bolzano στο διάστημα $[0,\pi]$, οπότε η εξίσωση $g(x)=0$ έχει τουλάχιστον μια ρίζα $\rho \in (0, \pi) \subset (0, +\infty)$. Επιπλέον η συνάρτηση $g$ είναι γνησίως αύξουσα στο διάστημα $[0, +\infty)$, οπότε η ρίζα $\rho$ είναι μοναδική στο διάστημα αυτό.
    Επειδή $g(-\rho)=(-\rho)^2-\sigma\upsilon\nu(-\rho)=\rho^2-\sigma\upsilon\nu \rho = g(\rho)=0$, άρα και το $-\rho$ είναι ρίζα της εξίσωσης $g(x)=0$. Επειδή $0 < \rho < \pi \Leftrightarrow -\pi < -\rho < 0$, η ρίζα $-\rho$ της εξίσωσης $g(x)=0$ βρίσκεται στο διάστημα $(-\pi,0)$.
    Επιπλέον η ρίζα $-\rho$ είναι μοναδική ρίζα της εξίσωσης $g(x)=0$ στο διάστημα $(-\infty, 0]$ αφού η συνάρτηση $g$ είναι γνησίως φθίνουσα στο διάστημα αυτό.
    Άρα η εξίσωση $g(x)=0 \Leftrightarrow x^2-\sigma\upsilon\nu x=0 \Leftrightarrow f^2(x)=5+\sigma\upsilon\nu x$ με $x \in \mathbb{R}$ έχει ακριβώς δύο ρίζες αντίθετες μεταξύ τους οι οποίες ανήκουν στο διάστημα $(-\pi, \pi)$.
\end{rlist}
\end{alist}